\let\textcircled=\pgftextcircled
\chapter{\'Etat de l'art}
\label{chap:etat-art}

\textit{\initial{C}e chapitre situe le travail dans le paysage de la r\'eponse
automatis\'ee aux incidents. Il examine successivement les briques de d\'etection,
les plateformes d'orchestration, les approches fond\'ees sur l'apprentissage, puis
d\'egage le verrou que la contribution entend lever.}

\mtcselectlanguage{french}
\minitoc
\newpage

\section{Cha\^ine de traitement d'un incident de s\'ecurit\'e}
\subsection{Cycle de vie normalis\'e}
% A COMPLETER : s'appuyer sur NIST SP 800-61r2 (Computer Security Incident
% Handling Guide) et ISO/IEC 27035. Les quatre phases : preparation ;
% detection et analyse ; confinement, eradication et retablissement ;
% activite post-incident.
Le cycle de vie retenu par la litt\'erature normative comporte quatre phases
[compl\'eter avec NIST SP 800-61r2 et ISO/IEC 27035]. La contribution de ce travail
porte sur la troisi\`eme --- le \textbf{confinement} --- et sur la boucle de retour
qui la relie \`a la quatri\`eme.

\subsection{Sources de d\'etection}
\begin{itemize}
  \item \textbf{\textsc{siem}} --- corr\'elation de journaux h\'et\'erog\`enes ;
  \item \textbf{\textsc{nids}} --- inspection du trafic r\'eseau ;
  \item \textbf{\textsc{edr}} --- t\'el\'em\'etrie des postes et serveurs ;
  \item \textbf{\textsc{ueba}} --- \'ecart \`a un profil comportemental de r\'ef\'erence.
\end{itemize}
[D\'evelopper : forces, limites, taux de faux positifs constat\'es dans la
litt\'erature.]

\section{Orchestration de la r\'eponse : les plateformes SOAR}
\subsection{Principe}
[Origine du terme, promesse : orchestration, automatisation, r\'eponse.]

\subsection{Limite structurelle}
La quasi-totalit\'e des plateformes de cette famille s'arr\^ete \`a la
\emph{recommandation}, ou n'automatise que des gestes sans effet sur la
production --- ouverture de ticket, enrichissement, notification. La raison n'est
pas technique : elle tient \`a l'absence de garantie sur la r\'eversibilit\'e des gestes
correctifs. Faute de pouvoir prouver qu'une action est annulable, l'\'editeur
reporte la responsabilit\'e sur l'analyste.

\begin{table}[H]
\centering
\caption{Comparaison des approches d'orchestration de la r\'eponse}
\label{tab:comparaison-soar}
\small
\begin{tabular}{@{}lcccc@{}}
\toprule
\textbf{Crit\`ere} & \textbf{SOAR classique} & \textbf{Playbooks} & \textbf{Approches ML} & \textbf{Ce travail} \\
 & & \textbf{scriptés} & & \\
\midrule
Action sans validation   & Non           & Partielle    & Variable & \textbf{Oui, born\'ee} \\
Garantie d'annulation    & Non trait\'ee  & Non trait\'ee & Non      & \textbf{Condition d'ex\'ecution} \\
Refus d'agir motiv\'e     & Non           & Non          & Non      & \textbf{Oui (ancrage)} \\
Explicabilit\'e           & \'Elev\'ee     & \'Elev\'ee    & Faible   & \'Elev\'ee \\
Boucle de contr\^ole      & Absente       & Absente      & Rare     & \textbf{Ferm\'ee} \\
Menace hors catalogue    & Silence       & Silence      & Extrapolation & \textbf{Repli par indicateurs} \\
Fonctionnement hors ligne& Non           & Oui          & Non      & \textbf{Oui} \\
\bottomrule
\end{tabular}
\end{table}

% A COMPLETER : citer et discuter au moins trois plateformes representatives,
% ainsi que deux travaux academiques sur la reponse automatisee.

\section{Approches fond\'ees sur l'apprentissage}
\subsection{Apport}
[D\'etection d'anomalies, r\'eduction du bruit, priorisation.]

\subsection{Obstacle \`a l'action autonome}
Trois obstacles interdisent de confier la \emph{d\'ecision d'agir} \`a un mod\`ele
appris dans le contexte vis\'e :

\begin{enumerate}
  \item \textbf{L'opposabilit\'e} --- une d\'ecision doit \^etre justifiable devant une
        autorit\'e ; un poids de r\'eseau ne constitue pas une justification ;
  \item \textbf{L'invention} --- un mod\`ele g\'en\'eratif produit une r\'eponse
        plausible m\^eme lorsqu'il ne dispose d'aucune donn\'ee fiable ;
  \item \textbf{La souverainet\'e} --- l'appel \`a un mod\`ele h\'eberg\'e hors du territoire
        contredit la contrainte de non-sortie des donn\'ees.
\end{enumerate}

C'est pourquoi, dans la conception retenue, le mod\`ele de langage n'est jamais
autoris\'e \`a produire un fait ni \`a choisir une action : il ne fait que reformuler
des faits d\'ej\`a \'etablis par le moteur d\'eterministe (voir
section~\ref{sec:assistant-conception}).

\section{R\'eversibilit\'e et s\^uret\'e de fonctionnement}
[Rapprocher des notions issues d'autres domaines : transactions et annulation en
bases de donn\'ees, patron \emph{Command}/\emph{Memento} en g\'enie logiciel,
\emph{circuit breaker} en syst\`emes r\'epartis, principe de r\'eversibilit\'e des
actions en s\^uret\'e des automatismes.]

\section{Synth\`ese et positionnement}
L'\'etat de l'art fait appara\^itre un verrou constant : l'automatisation de la
r\'eponse est trait\'ee comme un probl\`eme de \emph{confiance dans la d\'etection}, alors
qu'elle est d'abord un probl\`eme de \emph{co\^ut de l'erreur}. Le travail pr\'esent\'e ici
d\'eplace la question : plut\^ot que de chercher \`a rendre la d\'etection assez s\^ure pour
agir, il restreint l'action \`a ce dont l'erreur est rattrapable. La
r\'eversibilit\'e, jusque-l\`a trait\'ee comme une m\'etadonn\'ee documentaire, devient une
\textbf{condition op\'erationnelle} \'evalu\'ee avant chaque ex\'ecution.
